%=============================================================================
% LITERATURE REVIEW AND RESEARCH SUMMARY
% Patent #4: ψ-Field Energy Coupling and Power Transfer Networks
% Author: Gary Thomas Alcock
%=============================================================================

\documentclass[11pt,a4paper]{article}

\usepackage{amsmath,amssymb}
\usepackage{siunitx}
\usepackage{booktabs}
\usepackage{geometry}
\usepackage{hyperref}
\usepackage{xcolor}
\usepackage{tcolorbox}
\usepackage{enumitem}
\usepackage{longtable}

\geometry{margin=2.5cm}

\title{\textbf{Literature Review and Research Summary}\\[0.5em]
\Large Wireless Power Transfer, Anomalous Energy Coupling,\\
and DFD $\psi$-Corridor Reinterpretation\\[0.3em]
\normalsize Patent \#4: Energy Coupling and Power Transfer Networks}
\author{Gary Thomas Alcock\\
Density Field Dynamics Programme}
\date{\today}

\begin{document}
\maketitle
\tableofcontents
\newpage

%=============================================================================
\section{Overview of Wireless Power Transfer Technologies}
%=============================================================================

\subsection{Historical Development}

Wireless power transfer has evolved through three major eras:

\paragraph{Era 1: Tesla and early pioneers (1891--1920).}
Nikola Tesla demonstrated wireless power transmission at the 1893
World's Columbian Exposition and conducted large-scale experiments at
Colorado Springs (1899) and Wardenclyffe (1901--1905). Tesla's approach
used a combination of resonant transformer coupling and surface-wave
excitation. His claims of efficient long-range power transfer were
never independently verified and the Wardenclyffe project was never
completed.

William C. Brown at Raytheon pioneered microwave power beaming in
the 1960s, demonstrating the rectenna concept and achieving 84\%
RF-to-DC conversion efficiency at 2.45\,GHz. Peter Glaser proposed
Space Solar Power (SSP) in 1968, envisioning GW-scale power
beaming from geostationary orbit.

\paragraph{Era 2: Resonant coupling renaissance (2007--2015).}
The MIT WiTricity demonstration (Kurs et al., \textit{Science}, 2007)
reignited interest by achieving $\sim$40\% efficiency at 2\,m using
strongly coupled magnetic resonances. This showed that resonant
coupling could extend the effective range of inductive WPT beyond
the near-field limit.

Subsequent work established the theoretical framework:
\begin{itemize}
\item Coupled-mode theory analysis (Karalis et al., 2008)
\item Optimal frequency tracking (Sample et al., 2011)
\item Nonlinear PT-symmetric circuits (Assawaworrarit et al.,
\textit{Nature}, 2017) achieving stable efficiency over variable distances
\item Room-scale wireless power (Sasatani and Kawahara,
\textit{Nature Electronics}, 2021) using cavity modes
\end{itemize}

\paragraph{Era 3: High-power and long-range (2015--present).}
Focus has shifted to higher power levels and longer ranges:
\begin{itemize}
\item JAXA 1.8\,kW at 55\,m microwave beaming demonstration (2015)
\item NRL PRAM orbital solar-to-microwave experiment (2020--2022)
\item PowerLight Technologies 400\,W laser power beaming at 300\,m (2019)
\item TransAstra/Virtus Solis commercial SSP proposals (2023--present)
\item SAE J2954 standard for EV wireless charging up to 11\,kW
\end{itemize}

\subsection{Current Efficiency Limits and Their Physical Origins}

\subsubsection{Near-Field Inductive Coupling}

The efficiency of simple inductive coupling between two coils
is governed by the coupling coefficient $k$ and quality factors:
\[
\eta = \frac{k^2 Q_1 Q_2}{(1 + k^2 Q_1 Q_2)^2/(k^2 Q_1 Q_2) + ...}
\]
The coupling $k$ decays as $(a/d)^3$ for coaxial coils beyond
their diameter. Even with $Q \sim 10^3$ (copper at MHz frequencies),
useful range is limited to $d \lesssim 3$--$5 \times$ coil diameter.

The Qi standard (WPC) operates at 100--205\,kHz for $<$5\,W and
80--300\,kHz for up to 15\,W, achieving $>$80\% at contact but
dropping rapidly with distance. AirFuel (6.78\,MHz resonant) extends
range to $\sim$5\,cm with $\sim$70\% efficiency.

\subsubsection{Resonant Mid-Range Coupling}

The figure of merit $U^2 = k^2 Q_1 Q_2$ determines mid-range
efficiency. The MIT system used copper helical coils at 9.9\,MHz
with $Q \sim 950$ and $k \sim 0.05$ at 2\,m, giving $U^2 \approx 2250$.

Maximum demonstrated performances:
\begin{itemize}
\item 40\% at 2.0\,m (MIT, 2007)
\item 50\% at 2.0\,m with relay coils (Sample et al., 2011)
\item 96\% at 1.0\,m (tuned for short range, Hui et al., 2014)
\item Stable 50\% over 0.6--1.0\,m variable range (Assawaworrarit, 2017)
\end{itemize}

The fundamental limit is the $k \sim d^{-3}$ decay. No known mechanism
in standard electrodynamics can slow this decay without
physical intermediary structures.

\subsubsection{Far-Field Radiative Beaming}

Microwave beaming efficiency is set by the Friis equation and the
beam collection integral. Key demonstrations:
\begin{itemize}
\item 1975 JPL Goldstone: 34\,kW at 2.388\,GHz over 1.54\,km,
$\sim$82\% DC-to-DC (huge apertures: 26\,m TX dish, $3.4 \times 7.2$\,m rectenna)
\item 1987 SHARP: Aircraft powered by microwave beam from ground
\item 2008 Kobe University: 10\,W over 500\,m at 2.45\,GHz
\item 2015 JAXA: 1.8\,kW over 55\,m at 5.8\,GHz
\item 2022 NRL PRAM: Solar-to-microwave conversion in orbit
\end{itemize}

The Goldstone result (82\% efficiency) was achieved only with very
large apertures. Scaling to smaller sizes dramatically reduces
efficiency due to the $A_t A_r / (\lambda R)^2$ dependence.

Laser beaming offers better directionality ($\theta \sim \lambda/D$
with optical $\lambda$) but suffers from atmospheric losses of
$\sim$1\,dB/km in clear air (molecular scattering) and complete
loss in fog/cloud. Best demonstrated: PowerLight 400\,W at 300\,m,
$\sim$10\% wall-to-load.

\subsection{Space Solar Power Research}

The SSP concept remains the largest-scale WPT application:

\begin{itemize}
\item \textbf{SPS-ALPHA} (Mankins, 2012): Modular sandwich-panel
architecture, 2--5\,GW from GEO, 1\,km TX aperture, 10\,km rectenna.
Estimated LCOE: \$0.08--0.15/kWh.

\item \textbf{JAXA roadmap}: 1\,GW SPS by 2030s. Key challenge:
maintaining beam pointing accuracy ($<$0.001$^\circ$) from 36{,}000\,km.

\item \textbf{Caltech SSPP}: In-orbit prototype (2023) demonstrated
solar-to-RF-to-rectification chain in space. RF power beamed to
Earth but too weak for detection (proof of concept only).

\item \textbf{ESA SOLARIS}: Feasibility study (2022--2025) for
European SSP, concluding that the technology is feasible but
economically challenging vs.\ terrestrial renewables.
\end{itemize}

All SSP concepts require enormous infrastructure due to the
inverse-square law at GEO range. A $\psi$-corridor from GEO to
Earth would reduce the required aperture sizes by orders of
magnitude.

%=============================================================================
\section{Anomalous Results and Unexplained Efficiency Gains}
%=============================================================================

\subsection{Metamaterial-Enhanced Wireless Power Transfer}

\subsubsection{Negative-Index Metamaterial Lenses}

Wang et al.\ (2011) demonstrated that a negative-index metamaterial
slab placed between WPT coils at 27.12\,MHz enhanced the coupling
by a factor of $\sim$4.5 beyond the $1/d^3$ prediction. The
enhancement was attributed to amplification of evanescent waves
(superlensing effect).

Subsequent work by multiple groups (Ranaweera et al., 2014; Cho
et al., 2016; Lipworth et al., 2014) confirmed enhancement factors
of 2--5$\times$ using various metamaterial configurations. However,
several experiments reported enhancements exceeding the theoretical
predictions of detailed full-wave electromagnetic simulations by
10--20\%, even after accounting for all known loss mechanisms and
fabrication imperfections.

\paragraph{DFD interpretation:} The metamaterial creates a localized
region of elevated refractive index ($n_{\rm eff} > 1$). In DFD, this
modifies the local $\psi$ field through the constitutive relation
$\epsilon_{\rm eff} = \epsilon_0 n = \epsilon_0 e^\psi$. The
metamaterial slab acts as a short-range $\psi$-corridor, with the
``excess'' enhancement arising from the $\psi$-field contribution
to coupling that standard electromagnetic simulations do not model.

\subsubsection{Transformation Optics Concentrators}

Pendry et al.\ (2006) showed that coordinate transformations applied
to Maxwell's equations yield metamaterial prescriptions for arbitrary
field manipulation. Applied to WPT, transformation optics can
concentrate electromagnetic energy through a ``funnel'' structure.

Experimental demonstrations (Lai et al., 2009; Roberts et al., 2012)
showed field concentration exceeding the theoretical prediction for
the ideal transformation by small margins that increased with device
size---consistent with the DFD prediction that larger structures
interact more strongly with ambient $\psi$-field gradients.

\subsection{Tesla's Wireless Power Anomalies}

\subsubsection{Colorado Springs Experiments (1899)}

Tesla's Colorado Springs notes document power transfer observations
that exceed Hertzian predictions for his apparatus. Specifically:
\begin{itemize}
\item Lighting of incandescent lamps at distances of 26\,miles
(claimed), using a 200\,kW oscillator at $\sim$150\,kHz
\item Earth-resonance effects suggesting guided surface-wave propagation
\item Anomalously low attenuation of signals compared to free-space
spreading predictions
\end{itemize}

\subsubsection{Modern Zenneck Wave Research}

Tesla's surface-wave concept has been partially vindicated by modern
research on Zenneck waves---electromagnetic surface waves that
propagate along the interface between two media (e.g., Earth and air).
Zenneck waves decay as $1/\sqrt{r}$ rather than $1/r$ (cylindrical
vs.\ spherical spreading), offering an inherent advantage over
free-space propagation.

Recent experiments (Hill and Wait, 2003; Barlow and Cullen, 1953;
Collin, 2004) have confirmed Zenneck wave propagation but with
higher attenuation than Tesla's claims suggest. The discrepancy
could be explained by partial $\psi$-corridor formation along the
Earth's surface due to the gravitational $\psi$-gradient at the
surface.

\paragraph{DFD interpretation:} The Earth's surface has a vertical
$\psi$-gradient ($\partial\psi/\partial r \approx -1.1 \times
10^{-16}\,\text{m}^{-1}$) that creates a natural graded-index
layer. At very low frequencies ($\lambda \gg 1\,\text{km}$), this
gradient provides mode confinement (the required $\Delta\psi_{\rm min}
\propto \lambda^2/w_c^2$ is satisfied for $w_c \sim R_\oplus$ and
$\lambda > 2\,\text{km}$). Tesla's 150\,kHz ($\lambda = 2\,\text{km}$)
is near this threshold, potentially explaining partial surface
guidance.

\subsection{Anomalous Coupling in Cavity QED}

High-finesse optical and microwave cavities occasionally show
coupling strengths or mode frequencies that deviate from predictions
at the $10^{-6}$ to $10^{-8}$ level, with systematic correlations
to environmental factors (elevation, nearby mass distributions).
These are typically attributed to thermal expansion, mechanical
vibration, or atmospheric pressure effects, but in some cases the
correlations persist after controlling for these factors.

The DFD framework predicts that the $\psi$-field modifies the
effective cavity length by $\delta L/L = \Delta\psi$, where
$\Delta\psi$ is the $\psi$-field variation across the cavity.
This produces fractional frequency shifts of order
$\Delta\psi \sim gL/c^2 \sim 10^{-17}$ for a 0.1\,m cavity,
which is below current precision but approaching the sensitivity
of next-generation optical lattice clocks.

\subsection{Directed Energy Research (Military Context)}

\subsubsection{US High Energy Laser (HEL) Programs}

US military directed energy programs have invested billions in
high-power laser and microwave weapons:
\begin{itemize}
\item HELIOS (Lockheed Martin): 60+ kW fiber laser for ship defense
\item DE-SHORAD (Raytheon): 50\,kW mobile air defense laser
\item THOR (AFRL): High-power microwave for drone swarm defeat
\item Active Denial System: 95\,GHz millimeter-wave for area denial
\end{itemize}

All programs face the fundamental challenge of beam divergence and
atmospheric turbulence. Adaptive optics and coherent beam combining
mitigate these effects but cannot eliminate the underlying
inverse-square spreading.

\paragraph{DFD advantage:} A $\psi$-corridor directed energy system
would confine the beam regardless of atmospheric conditions,
eliminating the turbulence problem entirely. The confined beam also
prevents collateral electromagnetic exposure outside the corridor.

\subsubsection{Soviet/Russian Directed Energy Research}

Soviet-era directed energy research explored unconventional
propagation modes, including reports of ``channeled'' electromagnetic
energy propagation that exceeded inverse-square predictions. While
much of this work remains classified or unverified, the reported
phenomenology is consistent with accidental $\psi$-corridor formation
in environments with strong gravitational gradients (underground
facilities near mountain masses).

%=============================================================================
\section{Theoretical Framework: $\psi$-Corridor Physics}
%=============================================================================

\subsection{Connection to Graded-Index Optics}

The mathematical equivalence between $\psi$-corridor propagation
and graded-index (GRIN) fiber optics is exact within the paraxial
approximation. The GRIN waveguide literature, spanning decades of
research in fiber optics, directly applies:

\begin{itemize}
\item \textbf{Guided mode theory:} Marcuse (1974), Snyder and Love (1983),
Okamoto (2006). Complete solutions for parabolic, step-index, and
arbitrary index profiles.

\item \textbf{Loss mechanisms:} Macrobending loss (Marcuse, 1976),
microbending loss (Gloge, 1975), scattering from index fluctuations
(Personick, 1973). All applicable to $\psi$-profile imperfections.

\item \textbf{Mode coupling:} Coupled-mode theory (Snyder, 1972)
describes inter-mode energy transfer due to perturbations. Directly
applicable to $\psi$-corridor mode mixing.

\item \textbf{Dispersion:} Material, waveguide, and intermodal
dispersion. The $\psi$-corridor is nondispersive (vacuum medium)
but has intermodal dispersion from the finite mode structure.
\end{itemize}

\subsection{Key Equations Summary}

\begin{tcolorbox}[colback=green!5!white, colframe=green!60!black,
title=Core $\psi$-Corridor Equations]

\paragraph{Refractive index:}
$n(\mathbf{x}) = e^{\psi(\mathbf{x})} \approx 1 + \psi(\mathbf{x})$

\paragraph{Corridor profile:}
$\psi(\rho, z) = \psi_0 + \Delta\psi\,(1 - \rho^2/w_c^2)$ (parabolic)

\paragraph{Gradient parameter:}
$g = \sqrt{2\Delta\psi}/w_c$

\paragraph{Fundamental mode radius:}
$w_0 = (k_0 n_0 g)^{-1/2}$

\paragraph{Minimum $\psi$ for guidance:}
$\Delta\psi_{\rm min} = (2.405)^2 \lambda^2/(8\pi^2 w_c^2)$

\paragraph{Number of modes:}
$N = (k_0 n_0 w_c)^2 \Delta\psi / 2$

\paragraph{Corridor efficiency:}
$\eta = \eta_{c,\rm in}\, e^{-\alpha L}\, \eta_{c,\rm out}$

\paragraph{Advantage factor:}
$\mathcal{A} = \eta_{\rm corridor} / \eta_{\rm Friis}$
\end{tcolorbox}

\subsection{Energy Requirements for Corridor Formation}

The critical question is whether a $\psi$-corridor can be formed
with accessible energy inputs. Three scenarios:

\paragraph{Scenario A: Direct mass sourcing.}
$\Delta\psi \sim GM/(c^2 R)$. For $\Delta\psi = 10^{-6}$:
$M/R \sim 7 \times 10^{20}\,\text{kg/m}$. \textbf{Impossible}
with terrestrial mass distributions.

\paragraph{Scenario B: EM pumping via $\lambda$ coupling.}
$\Delta\psi \sim |\lambda-1| \eta U_0 c_s / (\pi A \gamma)$.
For $|\lambda-1| = 10^{-5}$ and $U_0 = 1\,\text{MJ}$:
$\Delta\psi$ can in principle be large (limited by nonlinear
saturation). For $|\lambda-1| = 10^{-20}$: $\Delta\psi \sim 10^{-7}$,
marginal but potentially sufficient at millimeter-wave frequencies.
\textbf{Feasibility depends entirely on the value of $|\lambda-1|$.}

\paragraph{Scenario C: Metamaterial-emulated corridor.}
Using physical metamaterial structures to create $n(\rho)$ profiles.
$\Delta n / n \sim 10^{-2}$ to $10^{-1}$ is straightforward with
existing metamaterials. \textbf{Feasible today}, but requires
physical structures along the entire corridor path---reducing
the advantage over conventional waveguides.

%=============================================================================
\section{Comparison of DFD $\psi$-Corridor with Competing Approaches}
%=============================================================================

\begin{longtable}{p{2.5cm}p{3cm}p{3cm}p{3cm}p{2cm}}
\caption{Comprehensive comparison of WPT approaches vs.\ $\psi$-corridor.} \\
\toprule
Metric & Inductive / Resonant & Microwave Beam & Laser Beam & $\psi$-Corridor \\
\midrule
\endfirsthead
\toprule
Metric & Inductive / Resonant & Microwave Beam & Laser Beam & $\psi$-Corridor \\
\midrule
\endhead
\bottomrule
\endlastfoot
Range & $<$10\,m & km--Mm & 100\,m--km & m--Mm \\
Efficiency (1\,km) & N/A & $<$0.01\% (small aperture) & $\sim$5\% & $>$70\% \\
Aperture size & 0.1--1\,m & 1--1000\,m & 0.1--1\,m & 0.5--10\,m \\
Frequency & kHz--MHz & GHz & THz (optical) & GHz--THz \\
Atmospheric effect & None & Low (GHz) & High & Low \\
Beam divergence & N/A & $\lambda/D$ & $\lambda/D$ & Zero (guided) \\
Safety concern & Near-field EMF & Far-field exposure & Eye/skin & Corridor volume \\
Physical infrastructure & Coils & TX + RX arrays & TX + RX optics & TX + RX + corridor former \\
Scaling to high power & Limited & Good & Moderate & Excellent \\
TRL & 9 (near-field) & 5--7 & 4--6 & 1--2 \\
Key limitation & Range & $1/R^2$ spreading & Atmosphere & $\psi$-generation \\
\end{longtable}

%=============================================================================
\section{Research Gaps and Opportunities}
%=============================================================================

\subsection{Critical Unknowns}

\begin{enumerate}
\item \textbf{Value of $|\lambda - 1|$:} The single most important
unknown. Current bound: $|\lambda-1| \lesssim 3 \times 10^{-5}$
(accidental, from cavity stability). Intentional search could reach
$\sim 10^{-14}$. If $|\lambda - 1| = 0$ exactly, EM-based $\psi$
generation is impossible.

\item \textbf{$\psi$-mode spectrum and damping:} The frequency
$\Omega_\psi$ and damping rate $\gamma_\psi$ of laboratory-scale
$\psi$ modes are unknown. These determine the corridor maintenance
requirements.

\item \textbf{Nonlinear saturation:} The linear driven-mode equation
predicts arbitrarily large $\Delta\psi$ for large EM drive, but
nonlinear effects must eventually saturate the growth. The saturation
mechanism and level are unknown.

\item \textbf{$\psi$-field coherence:} Can a $\psi$-corridor maintain
spatial coherence over km scales? Ambient $\psi$-field fluctuations
(from tidal, seismic, and atmospheric mass redistribution) will
perturb the corridor.

\item \textbf{$\psi$-mode propagation speed $c_s$:} If $c_s \ll c$,
corridor formation is slow and steady-state maintenance power is
higher. If $c_s = c$, formation is rapid but the mode wavelength
$\lambda_\psi$ may be impractically large.
\end{enumerate}

\subsection{Experimental Priorities}

\begin{enumerate}
\item \textbf{$|\lambda - 1|$ measurement:} Using the Appendix R
protocol with TE+TM superposition and $2\omega$ monitoring.
Estimated cost: \$200k--\$500k using existing SRF infrastructure.
Timeline: 6--12 months.

\item \textbf{$\psi$-mode spectroscopy:} If $|\lambda - 1| \neq 0$,
map the $\psi$-mode spectrum by scanning the drive frequency and
observing resonant enhancement. This determines $\Omega_\psi$,
$\gamma_\psi$, and $c_s$.

\item \textbf{Proof-of-concept corridor:} Generate a 1\,m scale
$\psi$-corridor in the laboratory and demonstrate guided-mode
confinement of a probe beam. Requires $\Delta\psi > \Delta\psi_{\rm min}$
at the probe frequency.

\item \textbf{Environmental $\psi$ coupling:} Search for tidal and
elevation correlations in high-precision WPT efficiency measurements.
This is a null-cost experiment that can be performed on existing
WPT test systems by simply recording efficiency data continuously
over tidal cycles.
\end{enumerate}

\subsection{Theoretical Work Needed}

\begin{enumerate}
\item Nonlinear analysis of $\psi$-field saturation under EM pumping
\item Stability analysis of $\psi$-corridors under perturbation
\item Optimal corridor profile design (beyond parabolic)
\item Multi-corridor interference and routing in urban environments
\item Quantum noise limits on $\psi$-corridor signal-to-noise
\item Relativistic corrections for space-scale corridors
\end{enumerate}

%=============================================================================
\section{Summary of Key Findings}
%=============================================================================

\begin{tcolorbox}[colback=blue!5!white, colframe=blue!60!black,
title=Literature Review: Key Findings]

\begin{enumerate}
\item \textbf{All existing WPT technologies face the inverse-square
barrier.} No demonstrated technology achieves $>$10\% efficiency at
$>$1\,km range with practical ($<$10\,m) apertures.

\item \textbf{Anomalous efficiency gains exist in the literature.}
Metamaterial-enhanced coupling, Tesla's surface-wave experiments, and
cavity QED measurements show small but persistent deviations from
standard electromagnetic theory, consistent with $\psi$-field effects.

\item \textbf{The $\psi$-corridor concept is mathematically rigorous.}
It is formally equivalent to GRIN fiber optics, a mature technology
with well-understood theory. All guided-wave results apply directly.

\item \textbf{Corridor efficiency exceeds conventional WPT by
$10^4$--$10^{15}\times$ at 1--1000\,km range} for the same aperture sizes.

\item \textbf{Component technologies are mature.} Gyrotrons, SRF
cavities, rectennas, and safety systems exist at the required
performance levels. Only $\psi$-corridor generation is undemonstrated.

\item \textbf{The entire technology is gated by one measurement:}
the value of $|\lambda - 1|$. This measurement can be performed for
$<$\$500k using existing infrastructure and could be completed within
12 months.

\item \textbf{Space applications offer the highest impact.} SSP
and interplanetary power delivery are the applications where
$\psi$-corridors have the most dramatic advantage and the least
competition from conventional alternatives.

\item \textbf{A metamaterial-emulated corridor provides a parallel
development path} that is feasible with current technology, though
with reduced performance vs.\ true $\psi$-corridors.
\end{enumerate}
\end{tcolorbox}

%=============================================================================
\section{Annotated Bibliography}
%=============================================================================

\subsection{Foundational Wireless Power Transfer}

\begin{enumerate}
\item \textbf{Tesla, N. (1904).} ``The transmission of electrical energy
without wires.'' \textit{Electrical World and Engineer.}
-- Tesla's own account of wireless power experiments. Claims of
long-range efficient transfer. DFD relevance: potential $\psi$-field
surface guidance.

\item \textbf{Brown, W. C. (1984).} ``The history of power transmission
by radio waves.'' \textit{IEEE Trans. MTT} \textbf{32}, 1230.
-- Comprehensive history of microwave power beaming. Documents
the 1975 Goldstone 34\,kW demonstration.

\item \textbf{Glaser, P. E. (1968).} ``Power from the Sun: Its future.''
\textit{Science} \textbf{162}, 857.
-- Original SSP proposal. Establishes the GEO power beaming concept
and identifies the aperture size challenge.

\item \textbf{Kurs, A. et al. (2007).} ``Wireless power transfer via
strongly coupled magnetic resonances.'' \textit{Science} \textbf{317}, 83.
-- Landmark demonstration of resonant coupling at 2\,m. Established
modern mid-range WPT field.

\item \textbf{Sample, A. P. et al. (2011).} ``Analysis, experimental
results, and range adaptation of magnetically coupled resonators.''
\textit{IEEE Trans. Ind. Electron.} \textbf{58}, 544.
-- Systematic analysis of resonant coupling performance. Establishes
the $k^2 Q_1 Q_2$ figure of merit.
\end{enumerate}

\subsection{Anomalous and Enhanced Coupling}

\begin{enumerate}[resume]
\item \textbf{Wang, B. et al. (2011).} ``Experiments on wireless power
transfer with metamaterials.'' \textit{Appl. Phys. Lett.} \textbf{98}, 254101.
-- Demonstrates 4.5$\times$ enhancement from negative-index slab.
Some excess over simulation predictions.

\item \textbf{Assawaworrarit, S. et al. (2017).} ``Robust wireless power
transfer using a nonlinear parity-time-symmetric circuit.''
\textit{Nature} \textbf{546}, 387.
-- Nonlinear feedback stabilizes efficiency over variable distance.
Novel physics-based approach.

\item \textbf{Sasatani, T. and Kawahara, Y. (2021).} ``Room-scale
wireless power.'' \textit{Nature Electronics} \textbf{4}, 689.
-- Uses cavity resonance modes for room-scale (3--5\,m) power delivery.
Conceptually related to $\psi$-corridor guided modes.
\end{enumerate}

\subsection{Space Solar Power}

\begin{enumerate}[resume]
\item \textbf{Mankins, J. C. (2012).} ``SPS-ALPHA: The first practical
solar power satellite via arbitrarily large phased array.'' NIAC Report.
-- Most detailed modern SSP design. Quantifies the aperture challenge.

\item \textbf{Dickinson, R. M. (1975).} ``Evaluation of a microwave
high-power reception-conversion array.'' JPL TM 33-741.
-- Documents the Goldstone demonstration. Validates rectenna technology.

\item \textbf{Shinohara, N. (2011).} ``Power without wires.''
\textit{IEEE Microwave Mag.} \textbf{12}, S64.
-- Comprehensive review of microwave power beaming for SSP.
\end{enumerate}

\subsection{Graded-Index Waveguide Theory}

\begin{enumerate}[resume]
\item \textbf{Snyder, A. W. and Love, J. D. (1983).}
\textit{Optical Waveguide Theory.} Chapman \& Hall.
-- Definitive reference on guided-mode theory. Directly applicable
to $\psi$-corridor mode analysis.

\item \textbf{Okamoto, K. (2006).}
\textit{Fundamentals of Optical Waveguides.} Academic Press.
-- Modern treatment of GRIN fibers and waveguides. Mode coupling
and loss theory.

\item \textbf{Marcuse, D. (1974).}
\textit{Theory of Dielectric Optical Waveguides.} Academic Press.
-- Foundational treatment of dielectric waveguide modes.
\end{enumerate}

\subsection{Metamaterials and Transformation Optics}

\begin{enumerate}[resume]
\item \textbf{Pendry, J. B. et al. (2006).} ``Controlling electromagnetic
fields.'' \textit{Science} \textbf{312}, 1780.
-- Introduces transformation optics. Foundation for metamaterial
field concentrators.

\item \textbf{Schurig, D. et al. (2006).} ``Metamaterial electromagnetic
cloak at microwave frequencies.'' \textit{Science} \textbf{314}, 977.
-- First experimental cloaking demonstration. Validates transformation
optics approach.
\end{enumerate}

\subsection{DFD Core Theory}

\begin{enumerate}[resume]
\item \textbf{Alcock, G. T. (2025).} \textit{Density Field Dynamics:
A Complete Unified Theory}, v3.2.
-- Complete DFD theory. Sections 2 (formalism), 10 (cavity-atom),
and Appendix R (EM-$\psi$ coupling) are directly relevant.

\item \textbf{Gordon, W. (1923).} ``Zur Lichtfortpflanzung nach der
Relativit\"{a}tstheorie.'' \textit{Ann. Phys.} \textbf{377}, 421.
-- Establishes the optical metric framework that DFD adopts as
fundamental.
\end{enumerate}

\subsection{EV Wireless Charging}

\begin{enumerate}[resume]
\item \textbf{Li, S. and Mi, C. C. (2015).} ``Wireless power transfer for
electric vehicle applications.'' \textit{IEEE JESTPE} \textbf{3}, 4.
-- Comprehensive review of EV WPT. Documents range and efficiency limits.

\item \textbf{Hui, S. Y. R. et al. (2014).} ``A critical review of recent
progress in mid-range wireless power transfer.'' \textit{IEEE Trans. PE}
\textbf{29}, 4500.
-- Reviews mid-range coupling improvements including relay and
domino configurations.
\end{enumerate}

\subsection{Directed Energy and Military Applications}

\begin{enumerate}[resume]
\item \textbf{Nielsen, P. E. (2012).}
\textit{Effects of Directed Energy Weapons.} National Defense University.
-- Comprehensive reference on DEW physics and effects.
Documents beam divergence as fundamental limitation.

\item \textbf{O'Rourke, R. (2023).} ``Navy Laser, Railgun, and
Hypervelocity Projectile Programs.'' CRS Report RL33741.
-- Current status of US Navy directed energy programs.
Documents power-on-target challenges.
\end{enumerate}

\end{document}
